\begin{exercise}{60}
Show that $J$ is not a compact simplicial set, i.e., that there are infinitely many
non-degenerate simplices in $J$.

\begin{proof}
By Definition~1.3.25, an $n$-simplex of $J$ is a tuple
$(a_0,\ldots,a_n)\in\{0,1\}^{n+1}$.
Faces delete entries and degeneracies repeat entries.
A tuple is therefore degenerate exactly when two adjacent
entries coincide. Hence the two alternating tuples are
nondegenerate in every dimension, giving infinitely many
nondegenerate simplices.

To prove noncompactness itself, use Exercise~\ref{ex:15} to write
$J=\operatorname*{colim}_n\mathrm{sk}_n(J)$.
The indexing poset $\mathbb N$ is filtered by
Exercise~\ref{ex:53}, since every finite subset has an upper bound.
If $J$ were compact, surjectivity of the canonical map on hom-sets
would make $\mathrm{id}_J$ factor as
\[
 J\xrightarrow{u}\mathrm{sk}_n(J)
 \xrightarrow{\ \iota_n\ }J
\]
for some $n$, where $\iota_n$ is the inclusion.
But by Exercise~\ref{ex:15}, $\mathrm{sk}_n(J)$ contains only
simplices whose nondegenerate term has dimension at most $n$.
It therefore misses the alternating $(n+1)$-simplex of $J$,
contradicting $\iota_nu=\mathrm{id}_J$.
\end{proof}
\end{exercise}
